% Options for packages loaded elsewhere
\PassOptionsToPackage{unicode}{hyperref}
\PassOptionsToPackage{hyphens}{url}
\PassOptionsToPackage{dvipsnames,svgnames,x11names}{xcolor}
\documentclass[
  a4paper,
]{article}
\usepackage{xcolor}
\usepackage[margin=20mm]{geometry}
\usepackage{amsmath,amssymb}
\setcounter{secnumdepth}{-\maxdimen} % remove section numbering
\usepackage{iftex}
\ifPDFTeX
  \usepackage[T1]{fontenc}
  \usepackage[utf8]{inputenc}
  \usepackage{textcomp} % provide euro and other symbols
\else % if luatex or xetex
  \usepackage{unicode-math} % this also loads fontspec
  \defaultfontfeatures{Scale=MatchLowercase}
  \defaultfontfeatures[\rmfamily]{Ligatures=TeX,Scale=1}
\fi
\usepackage{lmodern}
\ifPDFTeX\else
  % xetex/luatex font selection
  \setmainfont[]{DejaVu Serif}
  \setsansfont[]{DejaVu Sans}
  \setmonofont[]{DejaVu Sans Mono}
\fi
% Use upquote if available, for straight quotes in verbatim environments
\IfFileExists{upquote.sty}{\usepackage{upquote}}{}
\IfFileExists{microtype.sty}{% use microtype if available
  \usepackage[]{microtype}
  \UseMicrotypeSet[protrusion]{basicmath} % disable protrusion for tt fonts
}{}
\makeatletter
\@ifundefined{KOMAClassName}{% if non-KOMA class
  \IfFileExists{parskip.sty}{%
    \usepackage{parskip}
  }{% else
    \setlength{\parindent}{0pt}
    \setlength{\parskip}{6pt plus 2pt minus 1pt}}
}{% if KOMA class
  \KOMAoptions{parskip=half}}
\makeatother
\usepackage{longtable,booktabs,array}
\newcounter{none} % for unnumbered tables
\usepackage{calc} % for calculating minipage widths
% Correct order of tables after \paragraph or \subparagraph
\usepackage{etoolbox}
\makeatletter
\patchcmd\longtable{\par}{\if@noskipsec\mbox{}\fi\par}{}{}
\makeatother
% Allow footnotes in longtable head/foot
\IfFileExists{footnotehyper.sty}{\usepackage{footnotehyper}}{\usepackage{footnote}}
\makesavenoteenv{longtable}
\usepackage{graphicx}
\makeatletter
\newsavebox\pandoc@box
\newcommand*\pandocbounded[1]{% scales image to fit in text height/width
  \sbox\pandoc@box{#1}%
  \Gscale@div\@tempa{\textheight}{\dimexpr\ht\pandoc@box+\dp\pandoc@box\relax}%
  \Gscale@div\@tempb{\linewidth}{\wd\pandoc@box}%
  \ifdim\@tempb\p@<\@tempa\p@\let\@tempa\@tempb\fi% select the smaller of both
  \ifdim\@tempa\p@<\p@\scalebox{\@tempa}{\usebox\pandoc@box}%
  \else\usebox{\pandoc@box}%
  \fi%
}
% Set default figure placement to htbp
\def\fps@figure{htbp}
\makeatother
\setlength{\emergencystretch}{3em} % prevent overfull lines
\providecommand{\tightlist}{%
  \setlength{\itemsep}{0pt}\setlength{\parskip}{0pt}}
\usepackage{luatexja}
\usepackage{xurl}
\usepackage{float}
\usepackage{bookmark}
\IfFileExists{xurl.sty}{\usepackage{xurl}
\usepackage{float}}{} % add URL line breaks if available
\urlstyle{same}
\hypersetup{
  pdftitle={Electroweak Symmetry Breaking and Maximal Parity Violation in the 3D Triangle-Loop Product (Part IV)},
  pdfauthor={Noriaki Kihara (WF System Co., Ltd.) ORCID 0009-0004-6753-4020},
  colorlinks=true,
  linkcolor={Maroon},
  filecolor={Maroon},
  citecolor={Blue},
  urlcolor={Blue},
  pdfcreator={LaTeX via pandoc}}

\title{Electroweak Symmetry Breaking and Maximal Parity Violation in the
3D Triangle-Loop Product (Part IV)}
\usepackage{etoolbox}
\makeatletter
\providecommand{\subtitle}[1]{% add subtitle to \maketitle
  \apptocmd{\@title}{\par {\large #1 \par}}{}{}
}
\makeatother
\subtitle{― Exact 100\%:0\% Chiral Splitting in \(G_2 \supset SO(4)\),
Dynamical Higgs Mass Generation in \(\mathrm{Tri}(T)\),
\(W^\pm/Z^0/\gamma\) Spectrum, and Bare Weinberg Angle
\(\sin^2\theta_W^{\mathrm{bare}} = 3/8\) ―}
\author{Noriaki Kihara (WF System Co., Ltd.) ORCID 0009-0004-6753-4020}
\date{November 28, 2026 Version DOI 10.5281/zenodo.24626593}

\begin{document}
\maketitle

{
\setcounter{tocdepth}{2}
\tableofcontents
}
\textbf{Author:} Noriaki Kihara (WF System Co., Ltd.)\\
\textbf{ORCID:} 0009-0004-6753-4020\\
\textbf{Date:} November 28, 2026\\
\textbf{Version DOI:} 10.5281/zenodo.24626593\\
\textbf{Concept DOI:} 10.5281/zenodo.24626592\\
\textbf{Status:} Electroweak symmetry breaking and maximal parity
violation paper (\textbf{Part IV}) following Part I (Concept DOI:
\texttt{10.5281/zenodo.24622200}), Part II (Concept DOI:
\texttt{10.5281/zenodo.24622565}), and Part III (Concept DOI:
\texttt{10.5281/zenodo.24625962}). Starting from the \texttt{1+3+4}
split (pure 4-space \(SU(2)\)) of outer \(O_8^{\mathrm{ext}}\)
identified in Part III, we verify the electroweak sector
(\(SU(2)_L \times U(1)_Y \to U(1)_{\mathrm{em}}\), maximal parity
violation, dynamical \(W/Z\) mass generation, and bare Weinberg angle)
using deterministic machine-precision scripts
(\texttt{tri\_loop\_part4\_electroweak\_parity\_v0.py}, no random
numbers).

\begin{center}\rule{0.5\linewidth}{0.5pt}\end{center}

\subsection{Abstract}\label{abstract}

Parts I--III {[}P1--P3{]} established that the single background-free
real 3-tensor map
\(\mathrm{Tri}(T)_{ijk} = \sum_{p,q,r} T_{ipq} T_{qjr} T_{rpk}\) unifies
gravity (\texttt{1+3} past light cone and tidal splitting), QCD
(\(SU(3)_C\)), electromagnetism (\(U(1)_Q\)), one-generation CAR
fermions, and resolution-gauge readout, and identified a pure 3D
\(SU(2)\) subgroup on the extra 4-space \(\{e_4,e_5,e_6,e_7\}\) of outer
\(O_8^{\mathrm{ext}}\). To close one full generation of the Standard
Model, three questions remained: (1) why weak \(SU(2)_L\) couples
exclusively to one chirality (maximal parity violation), (2) how
hypercharge \(U(1)_Y\) commutes with all gauge and spatial symmetries,
and (3) whether dynamical \(\mathrm{Tri}(T)\) condensation generates
\(W/Z\) masses while leaving a massless residual symmetry.

In this Part IV, deterministic verification
(\texttt{tri\_loop\_part4\_electroweak\_parity\_v0.py}) establishes five
results:

\begin{enumerate}
\def\labelenumi{\arabic{enumi}.}
\tightlist
\item
  \textbf{Exact Algebraic Origin of Maximal Parity Violation (100\%:0\%
  Chiral Theorem in \(G_2\))}: Under
  \(SO(4) \cong SU(2)_{\mathrm{SD}} \times SU(2)_{\mathrm{ASD}}\) on
  \(\{e_4,e_5,e_6,e_7\}\), imposing the octonionic \(G_2\) (14D)
  derivation constraint forces the 3-space-invariant internal subgroup
  \(SU(2)_{\mathrm{pure4}}\) to occupy \textbf{self-dual (SD, +chiral)
  \texttt{1.000000000000} and anti-self-dual (ASD, −chiral)
  \texttt{0.000000000000}} (both \(G_2\) derivation residual and
  Lie-algebra residual \textbf{\texttt{0.000e+00}}). The opposite
  chirality (ASD) is locked 100\% to external 3D spatial rotations
  \(SO(3)\) and cannot act as an independent internal gauge force.
\item
  \textbf{Exact Closure of All 12 Standard Model Generators and Parity
  Selection Rules}: Defining the commuting hypercharge
  \(Y = I_8\otimes Q_{\mathrm{int}} - \frac{1}{2}P_4\otimes I_8\)
  (\(Q_{\mathrm{em}} = T_3^L + Y\)) on 64D
  \(O_8^{\mathrm{ext}}\otimes O_8^{\mathrm{int}}\), all commutators
  \([SO(3)_{\mathrm{ext}}, \mathrm{SM}_{12}]\) (36 pairs),
  \([SU(3)_C, SU(2)_L]\) (24 pairs), \([SU(3)_C, U(1)_Y]\) (8 pairs),
  and \([SU(2)_L, U(1)_Y]\) (3 pairs) vanish to
  \textbf{\texttt{0.000e+00}}. Under a \(W_1\) rotation
  (\(\theta=0.65\)), left-handed doublets (\(e_4\) sector) respond
  strongly (\texttt{0.323572}) while right-handed singlets (\(e_0, e_1\)
  sectors) give \textbf{\texttt{0.000e+00}}.
\item
  \textbf{Dynamical Higgs Condensation in \(\mathrm{Tri}(T)\): \(W\)
  Mass Generation vs.~Massless Residual}: On vacuum
  \(H = O_8\otimes M_8\), both \(SU(2)_L\) and residual 8D symmetries
  have breaking norm \(\le 3.8\times 10^{-14}\) (massless). Evolving a
  localized Higgs condensate P4\(_4\) at site \(A=4\)
  (\(\epsilon_H = 0.05 \to 0.22, \tau=40\)) under \(\mathrm{Tri}(T)\)
  generates 6-digit degenerate breaking norms \(\|W_a(U)\|_F =\)
  \textbf{\texttt{0.012593\ \$\textbackslash{}to\$\ 0.025105\ \$\textbackslash{}to\$\ 0.035323\ \$\textbackslash{}to\$\ 0.035517}}
  for \(W_1, W_2, W_3\), while the residual 8D symmetry fixing \(e_4\)
  remains strictly massless (\textbf{\texttt{3.82e-14}}).
\item
  \textbf{Bare Weinberg Angle \(\sin^2\theta_W^{\mathrm{bare}} = 3/8\)
  and 4×4 Electroweak Spectrum (\(\gamma, W^\pm, Z^0\))}: Without
  assuming \(SU(5)/SO(10)\), raw 64D trace ratios give
  \(\mathrm{Tr}(T_3^2)=8\), \(\mathrm{Tr}(Y^2)=40/3\),
  \(\mathrm{Tr}(Q_{\mathrm{em}}^2)=64/3\), yielding
  \textbf{\(\sin^2\theta_W^{\mathrm{bare}} = 3/8 = 0.375000\)} (matching
  the canonical GUT value to 6 digits). Substituting this parameter-free
  geometric ratio \(g'/g = \sqrt{3/5}\) into the 4×4 mass-squared matrix
  \(M^2\) of the complex Higgs doublet
  (\(T_3^L=-1/2, Y=+1/2 \Rightarrow Q_{\mathrm{em}}=0\)) produces
  eigenvalues
  \([m_\gamma^2, m_{W^+}^2, m_{W^-}^2, m_Z^2] = [\mathbf{-2.082\times 10^{-17}}, \mathbf{0.250000}, \mathbf{0.250000}, \mathbf{0.400000}]\)
  and
  \textbf{\(m_W^2/m_Z^2 = 5/8 = 0.625000 = \cos^2\theta_W^{\mathrm{bare}}\)}.
\item
  \textbf{Chiral Mirror Audit (\(M_8\) vs.~\(M_8^T\))}: Simultaneously
  mirroring background and condensate yields identical
  \(\|P\|=0.046825\) and \(\|W_1(U)\|_F=0.036160\), confirming that the
  Part I §8.3 \texttt{-1.41\%} chiral splitting arises purely from
  relative chirality between background and particle.
\end{enumerate}

\begin{center}\rule{0.5\linewidth}{0.5pt}\end{center}

\subsection{1. Introduction: Closing One Generation of the Standard
Model}\label{introduction-closing-one-generation-of-the-standard-model}

Part II {[}P2{]} established \(SU(3)_C \times U(1)_Q\) and CAR
\texttt{1:3:3:1} charges on 64D
\(O_8^{\mathrm{ext}}\otimes O_8^{\mathrm{int}}\), and Part III {[}P3{]}
isolated a 3-space-invariant pure \(SU(2)\) on the extra 4-space
\(\{e_4,e_5,e_6,e_7\}\) of \(O_8^{\mathrm{ext}}\). Closing
\(SU(3)_C \times SU(2)_L \times U(1)_Y\) requires answering three
questions: (1) why only one chirality of
\(SO(4) \cong SU(2)_L \times SU(2)_R\) survives as a weak internal
force, (2) what bare Weinberg angle emerges from the commuting
hypercharge \(Y\) with zero free parameters, and (3) whether
\(\mathrm{Tri}(T)\) dynamics itself spontaneously breaks \(SU(2)_L\) via
a localized condensate without an externally imposed Mexican-hat
potential. Every claim is tagged as \textbf{【Known】}, \textbf{【Part
I--III Established】}, \textbf{【Part IV Verified】}, or
\textbf{【Hypothesis / Unverified】}.

\begin{center}\rule{0.5\linewidth}{0.5pt}\end{center}

\subsection{\texorpdfstring{2. Exact 100\%:0\% Chiral Splitting in
\(G_2 \supset SO(4)\) 【Part IV Verified
EW-1】}{2. Exact 100\%:0\% Chiral Splitting in G\_2 \textbackslash supset SO(4) 【Part IV Verified EW-1】}}\label{exact-1000-chiral-splitting-in-g_2-supset-so4-part-iv-verified-ew-1}

The rotation algebra of 4-space \(\{e_4,e_5,e_6,e_7\}\) decomposes
orthogonally into self-dual (SD, +chiral) and anti-self-dual (ASD,
−chiral) subalgebras:
\(\mathfrak{so}(4) \cong \mathfrak{su}(2)_{\mathrm{SD}} \oplus \mathfrak{su}(2)_{\mathrm{ASD}}\)
【Known】. Projecting the 4×4 blocks of the \texttt{3+4}-preserving
subgroup
\(SO(4) = SU(2)_{\mathrm{pure4}} \times SU(2)_{\mathrm{space}} \subset G_2\)
onto SD and ASD gives (Fig. 1(a)):

{\def\LTcaptype{none} % do not increment counter
\begin{longtable}[]{@{}
  >{\raggedright\arraybackslash}p{(\linewidth - 6\tabcolsep) * \real{0.2222}}
  >{\raggedright\arraybackslash}p{(\linewidth - 6\tabcolsep) * \real{0.2222}}
  >{\centering\arraybackslash}p{(\linewidth - 6\tabcolsep) * \real{0.2778}}
  >{\centering\arraybackslash}p{(\linewidth - 6\tabcolsep) * \real{0.2778}}@{}}
\toprule\noalign{}
\begin{minipage}[b]{\linewidth}\raggedright
3D Subgroup of \(G_2\)
\end{minipage} & \begin{minipage}[b]{\linewidth}\raggedright
Action on 3-Space \(\{e_1,e_2,e_3\}\)
\end{minipage} & \begin{minipage}[b]{\linewidth}\centering
4-Space SD (+Chiral) Fraction
\end{minipage} & \begin{minipage}[b]{\linewidth}\centering
4-Space ASD (−Chiral) Fraction
\end{minipage} \\
\midrule\noalign{}
\endhead
\bottomrule\noalign{}
\endlastfoot
\textbf{\(SU(2)_{\mathrm{pure4}}\) (Internal Weak)} & Strictly zero (max
\texttt{8.327e-17}) & \textbf{\texttt{1.000000000000}} &
\textbf{\texttt{0.000000000000}} \\
\textbf{\(SU(2)_{\mathrm{space}}\) (Space-Locked)} & \(SO(3)\) spatial
rotation & \textbf{\texttt{0.000000000000}} &
\textbf{\texttt{1.000000000000}} \\
\end{longtable}
}

The standard \(SU(2)_L\) basis \((T_1, T_2, T_3)\) satisfies both the
\(G_2\) derivation condition and the \(\mathfrak{su}(2)\) Lie algebra to
\textbf{\texttt{0.000e+00}}. Thus the octonionic multiplication table
(\(G_2\)) strictly permits only one chirality (SD) to act as a
3-space-commuting internal gauge force while locking the opposite
chirality (ASD) 100\% to spatial rotations \(SO(3)\)---providing a
direct algebraic origin for maximal parity violation 【Part IV
Verified】+【Hypothesis】.

\begin{center}\rule{0.5\linewidth}{0.5pt}\end{center}

\subsection{3. Exact Commutativity of All 12 Standard Model Generators
【Part IV Verified
EW-2】}\label{exact-commutativity-of-all-12-standard-model-generators-part-iv-verified-ew-2}

On 64D \(O_8^{\mathrm{ext}}\otimes O_8^{\mathrm{int}}\), using the
4-space projector \(P_4 = \mathrm{diag}(0,0,0,0,1,1,1,1)\), we define
the commuting hypercharge
\[Y \equiv I_8 \otimes Q_{\mathrm{int}} - \frac{1}{2}\, P_4 \otimes I_8, \qquad Q_{\mathrm{em}} = T_3^L + Y. \tag{1}\]
All commutators vanish to machine zero: -
\([SO(3)_{\mathrm{ext}},\; SU(3)_C \times SU(2)_L \times U(1)_Y]\) (36
pairs): \textbf{\texttt{0.000e+00}} - \([SU(3)_C,\; SU(2)_L]\) (24
pairs), \([SU(3)_C,\; U(1)_Y]\) (8 pairs), \([SU(2)_L,\; U(1)_Y]\) (3
pairs): \textbf{\texttt{0.000e+00}}

Under \(U_W = \exp(0.65\, T_1^L\otimes I_8)\), left-handed doublet
states (\(e_4\) sector) give
\(\|U_W\psi_L - \psi_L\| = \mathbf{0.323572}\), whereas right-handed
singlet states (\(e_0, e_1\) sectors) give \textbf{\texttt{0.000e+00}}.

\begin{center}\rule{0.5\linewidth}{0.5pt}\end{center}

\subsection{\texorpdfstring{4. Dynamical Higgs Condensation,
\(W/Z/\gamma\) Spectrum, and
\(\sin^2\theta_W^{\mathrm{bare}} = 3/8\)}{4. Dynamical Higgs Condensation, W/Z/\textbackslash gamma Spectrum, and \textbackslash sin\^{}2\textbackslash theta\_W\^{}\{\textbackslash mathrm\{bare\}\} = 3/8}}\label{dynamical-higgs-condensation-wzgamma-spectrum-and-sin2theta_wmathrmbare-38}

\subsubsection{\texorpdfstring{4.1 Dynamical \(W_1, W_2, W_3\) Mass
Generation vs.~Massless Residual in \(\mathrm{Tri}(T)\) 【Part IV
Verified
EW-3(a,b)】}{4.1 Dynamical W\_1, W\_2, W\_3 Mass Generation vs.~Massless Residual in \textbackslash mathrm\{Tri\}(T) 【Part IV Verified EW-3(a,b)】}}\label{dynamical-w_1-w_2-w_3-mass-generation-vs.-massless-residual-in-mathrmtrit-part-iv-verified-ew-3ab}

Without introducing any external potential, we place a localized
condensate P4\(_4\) at site \(A=4\) on vacuum \(H = O_8\otimes M_8\),
evolve \(\mathrm{Tri}(T)\) for \(\tau=40\) steps, and measure the
symmetry-breaking order parameter \(\|G(U_{\mathrm{higgs}})\|_F\) (Table
1, Fig. 1(b)).

{\def\LTcaptype{none} % do not increment counter
\begin{longtable}[]{@{}
  >{\centering\arraybackslash}p{(\linewidth - 10\tabcolsep) * \real{0.1667}}
  >{\centering\arraybackslash}p{(\linewidth - 10\tabcolsep) * \real{0.1667}}
  >{\centering\arraybackslash}p{(\linewidth - 10\tabcolsep) * \real{0.1667}}
  >{\centering\arraybackslash}p{(\linewidth - 10\tabcolsep) * \real{0.1667}}
  >{\centering\arraybackslash}p{(\linewidth - 10\tabcolsep) * \real{0.1667}}
  >{\centering\arraybackslash}p{(\linewidth - 10\tabcolsep) * \real{0.1667}}@{}}
\toprule\noalign{}
\begin{minipage}[b]{\linewidth}\centering
Condensate Amplitude \(\epsilon_H\)
\end{minipage} & \begin{minipage}[b]{\linewidth}\centering
\(\|W_1(U)\|_F\)
\end{minipage} & \begin{minipage}[b]{\linewidth}\centering
\(\|W_2(U)\|_F\)
\end{minipage} & \begin{minipage}[b]{\linewidth}\centering
\(\|W_3(U)\|_F\)
\end{minipage} & \begin{minipage}[b]{\linewidth}\centering
Residual 8D Max
\end{minipage} & \begin{minipage}[b]{\linewidth}\centering
\(E_{\mathrm{probe}}\)
\end{minipage} \\
\midrule\noalign{}
\endhead
\bottomrule\noalign{}
\endlastfoot
\texttt{0.00} (Vacuum \(H\)) & \texttt{0.000000} & \texttt{0.000000} &
\texttt{0.000000} & \texttt{3.802e-14} & \texttt{0.000000} \\
\texttt{0.05} & \textbf{\texttt{0.012593}} & \textbf{\texttt{0.012593}}
& \textbf{\texttt{0.012593}} & \texttt{3.821e-14} & \texttt{0.000294} \\
\texttt{0.10} & \textbf{\texttt{0.025105}} & \textbf{\texttt{0.025105}}
& \textbf{\texttt{0.025105}} & \texttt{3.821e-14} & \texttt{0.000964} \\
\texttt{0.16} & \textbf{\texttt{0.035323}} & \textbf{\texttt{0.035323}}
& \textbf{\texttt{0.035323}} & \texttt{3.828e-14} & \texttt{0.001510} \\
\texttt{0.22} & \textbf{\texttt{0.035517}} & \textbf{\texttt{0.035517}}
& \textbf{\texttt{0.035517}} & \texttt{3.821e-14} & \texttt{0.001055} \\
\end{longtable}
}

All three \(SU(2)_L\) generators acquire identical 6-digit degenerate
breaking norms upon condensation while the residual 8D symmetry fixing
\(e_4\) remains strictly massless (\texttt{3.82e-14}).

\begin{figure}[H]
\centering
\includegraphics[width=0.96\linewidth,height=\textheight,keepaspectratio,alt={Part IV Verification (1): (a) Exact 100\%:0\% SD/ASD chiral splitting in G\_2 \textbackslash supset SO(4) (algebraic origin of maximal parity violation); (b) Dynamical \textbackslash mathrm\{Tri\}(T) Higgs condensate P4\_4 generating degenerate W\_1, W\_2, W\_3 breaking norms vs.~massless residual 8D symmetry.}]{figs_en/fig_p4_parity_and_higgs_mass_en.png}
\caption{Part IV Verification (1): (a) Exact 100\%:0\% SD/ASD chiral
splitting in \(G_2 \supset SO(4)\) (algebraic origin of maximal parity
violation); (b) Dynamical \(\mathrm{Tri}(T)\) Higgs condensate P4\(_4\)
generating degenerate \(W_1, W_2, W_3\) breaking norms vs.~massless
residual 8D symmetry.}
\end{figure}

\subsubsection{\texorpdfstring{4.2 Bare Weinberg Angle
\(\sin^2\theta_W^{\mathrm{bare}} = 3/8\) and 4×4 Electroweak Mass
Spectrum 【Part IV Verified EW-3(c),
EW-5】}{4.2 Bare Weinberg Angle \textbackslash sin\^{}2\textbackslash theta\_W\^{}\{\textbackslash mathrm\{bare\}\} = 3/8 and 4×4 Electroweak Mass Spectrum 【Part IV Verified EW-3(c), EW-5】}}\label{bare-weinberg-angle-sin2theta_wmathrmbare-38-and-44-electroweak-mass-spectrum-part-iv-verified-ew-3c-ew-5}

Evaluating the trace-squared norms over 64D
\(O_8^{\mathrm{ext}}\otimes O_8^{\mathrm{int}}\) yields (Fig. 2(b)):
\[\mathrm{Tr}(T_3^2) = 8.000000,\qquad \mathrm{Tr}(Y^2) = \frac{40}{3} = 13.333333,\qquad \mathrm{Tr}(Q_{\mathrm{em}}^2) = \frac{64}{3} = 21.333333, \tag{2}\]
which directly gives the parameter-free geometric mixing angle {[}E4{]}:
\[\sin^2\theta_W^{\mathrm{bare}} = \frac{\mathrm{Tr}(T_3^2)}{\mathrm{Tr}(Q_{\mathrm{em}}^2)} = \frac{3}{8} = \mathbf{0.375000}, \qquad \frac{g'}{g} = \sqrt{\frac{3}{5}}. \tag{3}\]
Diagonalizing the 4×4 mass-squared matrix
\(M_{ab}^2 = \langle G_a\phi_H, G_b\phi_H\rangle\) for the neutral
component of the complex Higgs doublet
(\(T_3^L=-1/2, Y=+1/2 \Rightarrow Q_{\mathrm{em}}=0\)) with
\(g'/g=\sqrt{3/5}\) yields (Fig. 2(a)):
\[[m_\gamma^2,\; m_{W^+}^2,\; m_{W^-}^2,\; m_Z^2] = [\mathbf{-2.082\times 10^{-17}},\; \mathbf{0.250000},\; \mathbf{0.250000},\; \mathbf{0.400000}], \qquad \frac{m_W^2}{m_Z^2} = \mathbf{0.625000} = \frac{5}{8} = \cos^2\theta_W^{\mathrm{bare}}. \tag{4}\]

\begin{figure}[H]
\centering
\includegraphics[width=0.96\linewidth,height=\textheight,keepaspectratio,alt={Part IV Verification (2): (a) 4×4 electroweak mass-squared spectrum showing exact massless photon \textbackslash gamma, degenerate W\^{}\textbackslash pm, and m\_W\^{}2/m\_Z\^{}2 = 5/8; (b) Bare Weinberg angle \textbackslash sin\^{}2\textbackslash theta\_W\^{}\{\textbackslash mathrm\{bare\}\} = 3/8 from raw 64D trace ratios.}]{figs_en/fig_p4_ew_spectrum_and_weinberg_en.png}
\caption{Part IV Verification (2): (a) 4×4 electroweak mass-squared
spectrum showing exact massless photon \(\gamma\), degenerate \(W^\pm\),
and \(m_W^2/m_Z^2 = 5/8\); (b) Bare Weinberg angle
\(\sin^2\theta_W^{\mathrm{bare}} = 3/8\) from raw 64D trace ratios.}
\end{figure}

\subsubsection{\texorpdfstring{4.3 Chiral Mirror Audit (\(M_8\)
vs.~\(M_8^T\)) 【Part IV Verified
EW-4】}{4.3 Chiral Mirror Audit (M\_8 vs.~M\_8\^{}T) 【Part IV Verified EW-4】}}\label{chiral-mirror-audit-m_8-vs.-m_8t-part-iv-verified-ew-4}

Simultaneously mirroring the inner time crystal (\(M_8 \to M_8^T\)) for
both background and condensate produces identical \(\|P\| = 0.046825\)
and \(\|W_1(U)\|_F = 0.036160\), confirming that the Part I §8.3
\texttt{-1.41\%} asymmetry is a genuine relative-chirality effect
between background and particle.

\begin{center}\rule{0.5\linewidth}{0.5pt}\end{center}

\subsection{5. Clean Separation of Remaining Open
Problems}\label{clean-separation-of-remaining-open-problems}

The following items remain \textbf{【Hypothesis / Unverified】} and are
explicitly deferred to Parts V--VII: 1. \textbf{Renormalization-group
running (\(\sin^2\theta_W: 3/8 \to 0.231\))}: Deriving the scale
evolution from the bare geometric value \(3/8\) to the low-energy
measured value \(\approx 0.231\). 2. \textbf{Multi-shell
\(O_8^{\otimes K}\) continuum limit (planned Part V)}: \(1/r\) gravity,
\(1/r^2\) Coulomb force, light-cone speed \(c\), Einstein's equations,
and cumulative gravitational time dilation \(g_{00}(r)\). 3.
\textbf{Dynamical scattering \& multi-body quantum correlations (planned
Part VI)}: Two-particle scattering and dynamical CHSH Bell inequality
violation. 4. \textbf{Three generations, precision mass ratios, and
coupling constants (planned Part VII)}: \(J_3(\mathbb{O})\) /
\(O_8^{\otimes K}\) three-generation structure, exact mass ratios, and
first-principles derivation of \(\alpha \approx 1/137.036\) and
\(\alpha_s\).

\begin{center}\rule{0.5\linewidth}{0.5pt}\end{center}

\subsection{6. Reproducibility and Published
Files}\label{reproducibility-and-published-files}

\begin{itemize}
\tightlist
\item
  Verification script:
  \texttt{tri\_loop\_part4\_electroweak\_parity\_v0.py} (deterministic,
  \textasciitilde15 s)
\item
  Numerical log:
  \texttt{results/tri\_loop\_part4\_electroweak\_parity\_v0.txt}
\item
  Figure script: \texttt{plot\_part4\_figures\_v0.py}
\item
  Figures:
  \texttt{results/fig\_p4\_parity\_and\_higgs\_mass\_\{ja,en\}.png},
  \texttt{results/fig\_p4\_ew\_spectrum\_and\_weinberg\_\{ja,en\}.png}
\end{itemize}

\begin{center}\rule{0.5\linewidth}{0.5pt}\end{center}

\subsection{References}\label{references}

\subsubsection{A. Self-Citations}\label{a.-self-citations}

\begin{itemize}
\tightlist
\item
  \textbf{{[}P1{]}} N. Kihara, \emph{Spontaneous Emergence of Quantum,
  Gauge, Gravitational, and Fermionic Structures from the 3D
  Triangle-Loop Product (Unified Full Paper, Part I)}, Zenodo
  (2026-11-28), Concept DOI: \texttt{10.5281/zenodo.24622200}.
\item
  \textbf{{[}P2{]}} N. Kihara, \emph{Unified Connection Hypothesis for
  Gravity, Electromagnetism, QCD, and Dirac Matter in the 3D
  Triangle-Loop Product (Part II)}, Zenodo (2026-11-28), Concept DOI:
  \texttt{10.5281/zenodo.24622565}.
\item
  \textbf{{[}P3{]}} N. Kihara, \emph{Spacetime and Internal-Space
  Readout Connection Hypothesis in the 3D Triangle-Loop Product (Part
  III)}, Zenodo (2026-11-28), Concept DOI:
  \texttt{10.5281/zenodo.24625962}.
\end{itemize}

\subsubsection{B. External Primary
Literature}\label{b.-external-primary-literature}

\begin{itemize}
\tightlist
\item
  \textbf{{[}E1{]}} S. Weinberg, \emph{A Model of Leptons}, Phys.
  Rev.~Lett. \textbf{19}, 1264--1266 (1967); A. Salam, in
  \emph{Elementary Particle Theory}, ed.~N. Svartholm (1968).
\item
  \textbf{{[}E2{]}} P. W. Higgs, \emph{Broken Symmetries and the Masses
  of Gauge Bosons}, Phys. Rev.~Lett. \textbf{13}, 508--509 (1964).
\item
  \textbf{{[}E3{]}} C. Furey,
  \emph{\(SU(3)_C \times SU(2)_L \times U(1)_Y (\times U(1)_X)\) as a
  symmetry of division algebraic ladder operators}, Eur. Phys. J. C
  \textbf{78}, 375 (2018), arXiv:1806.00612.
\item
  \textbf{{[}E4{]}} H. Georgi and S. L. Glashow, \emph{Unity of All
  Elementary-Particle Forces}, Phys. Rev.~Lett. \textbf{32}, 438--441
  (1974).
\item
  \textbf{{[}E5{]}} J. C. Baez, \emph{The Octonions}, Bull. Amer. Math.
  Soc. \textbf{39}(2), 145--205 (2002), arXiv:math/0105155.
\end{itemize}

\end{document}
